% ==============================================================================
% sec:teorie --- Přehled současného poznání
% ==============================================================================
% Zdroje: Hala_SystemySocialnihoDialogu2013, CMKOS_ZpravaKV2025, CMKOS_PrinosKV2025,
%         Armstrong_HRM2006, Dvorakova_AEP_SustainHRM2023, Dvorakova_DemografickeTrendy2025,
%         ETUI_RoadTo80_Muller_2025, Schulten_WSI_Tarifstaerkung_2024, etui_cba,
%         OECD_GoodJobsForAll_2018, OECD_UnionEmployerCoverage_2025, Pokorny_KVsolidarita2024,
%         Eurofound_KeyDimensionsIR_2023, Eurofound_CBBeyondPay_2024, eurofound_ca_pay_equity_2026,
%         EF25044_PayEquity_2026, WojtczukTurek_SustainHRM_JobSatisfaction2024, PAQ_Svarcsystem,
%         zakon_zp_2006, MPSV_AkcniPlanKV2025, MPSV_JMHZ_2026, Koncepce_DalsihoVzdelavani_SP,
%         MZCR_AkcniPlanDusevniZdravi2020, MPSV_ZakonOdmenovani_2026,
%         PennyMarket_KS2026, KS_CeskaPosta_2025, KSVS_OSPZV_ASO_2026, KSVS_OSPPP_2026
% ==============================================================================

Použitá literatura o~sociálním dialogu sestává z~právního výkladu, komparativního výzkumu pracovních vztahů, manažerského pojetí \acgen{HRM} a~empirického sledování \acgen{KV}. Závěrečné práce právnické, například text Jiřího Svobody \emph{Sociální dialog} nebo práce Ivy Paštikové \emph{Kolektivní vyjednávání a~kolektivní spory}, sledují především subjekty, proces uzavírání \acgen{KS}, kolektivní spory, stávku, výluku a~judikaturu. Závěrečné práce z~oblasti personálního řízení se naopak soustředí na~mzdovou politiku, benefity, motivaci, transparentnost výkonových složek, koncept \textit{Total Rewards} a~značku zaměstnavatele. Tato práce oba proudy doplňuje otázkou, zda sociální dialog dokáže přerozdělit náklady strukturálních změn trhu práce místo toho, aby byly přenášeny na~jednotlivé pracovníky, a~docílit tak udržitelného rozvoje hospodářství.\par


Komparativní mezinárodní argumentace je v~této práci systematicky soustředěna do~sekce \enquote{Evropské modely trhu práce}, kde jsou jednotlivé mechanismy rozpracovány po zemích a~vztahovány ke~slabinám českého uspořádání.\par

\paragraph{Manažerská literatura k~\acs{HRM}} sleduje strategii organizace, motivaci, odměňování, rozvoj pracovníků a~firemní kulturu. Armstrongovo pojetí \acgen{HRM} je pro tuto práci důležité zejména rolí \textit{employee advocate} v~Ulrichově modelu HR partnership. Personální útvar má být hlasem zaměstnanců uvnitř organizace, současně však zůstává součástí hierarchie zaměstnavatele. Zde vzniká limit interního zastoupení -- může tlumit konflikty a~přenášet informace, ale z~principu nedisponuje nezávislostí vůči vedení.\par

\Ac{KV} tento střet řeší jiným mechanismem. Nezávislá odborová organizace odděluje zaměstnanecký hlas od~personální linie zaměstnavatele, přičemž zákoník práce poskytuje odborovým funkcionářům zvláštní ochranu. Proto nejde o~konkurenci moderního \acgen{HRM}, nýbrž o~jeho doplněk. HR business partner může zprostředkovat organizační změnu, formulovat politiku odměňování nebo podporovat vzdělávání -- \ac{KV} převádí vybrané standardy do~vymahatelných pravidel. Prakticky se to týká mzdy, hodnocení výkonu, diverzity, \acgen{BOZP}, kvality pracovního života i~podnikového vzdělávání, tedy oblastí, které se v~malých a~středních podnicích často neřídí robustním HR aparátem.~\cite{Armstrong_HRM2006}\par

Současný výzkum udržitelného \aclgen{HRM} tento vztah dále zpřesňuje. Mezikulturní studie 54~zemí ukazuje, že udržitelná personální politika souvisí s~vyšší pracovní spokojeností a~že její účinek je silnější u~zaměstnanců s~nižší organizační identifikací. Jinými slovy, dobře nastavená pravidla a~participace mohou částečně nahrazovat slabší loajalitu k~organizaci. Sociální dialog v~tomto pohledu není pouze odborovou procedurou, ale jedním z~kanálů, jímž se sociální pilíř udržitelného řízení lidí stává kvantifikovatelným a~kontrolovatelným.\cite{WojtczukTurek_SustainHRM_JobSatisfaction2024}\par

\paragraph{Česká empirická základna} poskytuje praktické části oporu především každoroční zprávou \acgen{ČMKOS}. V~roce 2025 bylo v~podnikatelské sféře evidováno \num{3378} \ac{PKS} u~\num{4061} zaměstnavatelů, které pokrývaly přibližně \SI{1382019}{\person}. \Ac{KSVS} pokrývaly \SI{440410}{\person}, z~toho \SI{160338}{\person} prostřednictvím rozšíření závaznosti. Význam extenze je tedy zřejmý: počet smluv vyššího stupně klesá, ale počet pokrytých zaměstnavatelů roste právě tam, kde se závaznost rozšiřuje mimo členskou základnu zaměstnavatelského svazu. Tato zpráva dokládá, že uzavření \ac{KS} přináší zaměstnancům lepší podmínky, tedy vyšší mzdy a~nižší odpracovanou dobu v~podnicích s~kolektivní smlouvou oproti podnikům bez ní.~\cite{CMKOS_ZpravaKV2025}\par

Tato evidence je však zároveň neúplná. Česká data o~pokrytí \acins{KS} dlouho stojí na~kombinaci odborových šetření, \acloc{ISPV} a~odhadů, které postačují pro určení, že \ac{ČR} nepřekračuje evropský práh \SI{80}{\percent}, ale nestačí pro přesné cílení politiky \acgen{KV}. Akční plán \acgen{MPSV} proto staví první okruh opatření na~zlepšení datové základny, zejména na~\ac{JMHZ}, které má od~roku 2026 zachytit typ závazné kolektivní smlouvy či dohody, odvětví podle \acgen{NACE} a~kombinaci podnikové, vyšší a~rozšířené smlouvy u~všech zaměstnavatelů. První robustní sektorová data mají být dostupná až v~roce 2027, pročež současná praktická část musí pracovat s~přechodným stavem dat.~\cite{MPSV_AkcniPlanKV2025}~\cite{MPSV_JMHZ_2026}\par

\paragraph{Adaptace trhu práce a~\acloc{APZ}} posouvá sociální dialog od~mzdové redistribuce k~otázce zaměstnatelnosti. Český trh práce se vyznačuje nízkou nezaměstnaností, ale současně trpí nesouladem dovedností, omezenou mobilitou a~nedostatečnou prevencí ztráty zaměstnání u~starších pracovníků. \Ac{APZ} je v~českém prostředí převážně reaktivní, zatímco demografické stárnutí vyžaduje preventivní age-management, celoživotní vzdělávání, ergonomické standardy a~práci s~psychosociálními riziky ještě před ztrátou pracovního místa.~\cite{Dvorakova_AEP_SustainHRM2023}~\cite{Dvorakova_DemografickeTrendy2025}\par

Sociální partneři mají v~této oblasti dvojí roli. Zaprvé znají odvětvovou poptávku po~dovednostech lépe než centrální správa a~zadruhé mohou standardy vzdělávání, rekvalifikací a~placeného volna zakotvit v~\acploc{KS}. Návrh koncepce dalšího vzdělávání zpracovaný \acs{ČMKOS} a~\acs{SPČR} proto představuje důležitý most mezi veřejnou \acins{APZ}, podnikovým vzděláváním a~odvětvovým vyjednáváním. Ukazuje zároveň slabinu české praxe: konkrétní programy odborného rozvoje se v~podnikových smlouvách objevují jen zřídka, ačkoliv právě v~této oblasti lze kolektivní smlouvu chápat jako nástroj adaptace na~digitalizaci, automatizaci a~stárnutí pracovní síly.~\cite{Koncepce_DalsihoVzdelavani_SP}~\cite{InovacniDoporuceni_ZahraniceKV}\par

\paragraph{Nové trendy} rozšiřují pojetí \acgen{KV} za hranici mzdové otázky. Analýza \acpgen{KS} v~nízkomzdových odvětvích ukazuje, že smlouvy sice vždy regulují mzdy, příplatky a~pracovní dobu, avšak rozdíly mezi zeměmi spočívají v~míře detailu a~v~tom, zda se smlouvy zabývají také vzděláváním, přechody mezi pracemi, \acins{BOZP}, work-life balance, zelenou a~digitální transformací nebo psychosociálním zdravím. \ac{ČR} je ve vzorku zařazena mezi země, jejichž odvětvové smlouvy jsou spíše obecné. Pro rozvoj kolektivního vyjednávání je to důležité -- pouhé zvýšení počtu smluv nemusí stačit, pokud smlouvy neobsahují standardy, které reagují na~nová rizika na trhu práce.~\cite{Eurofound_CBBeyondPay_2024}\par

Zvláštním případem nové agendy je mzdová transparentnost. Směrnice o~transparentním odměňování přesouvá důkazní i~datovou zátěž směrem k~zaměstnavatelům a~posiluje roli zástupců zaměstnanců při společném posouzení odměňování. Kolektivní smlouvy mohou v~tomto rámci vytvářet genderově neutrální klasifikace prací, tarifní struktury a~veřejně čitelné kategorie pracovníků.~\cite{eu_directive_2022_2041}~\cite{eurofound_ca_pay_equity_2026}\par


Kvalita pracovního prostředí mimoto zahrnuje psychosociální rizika -- nové technologie, časový tlak, ztrátu kontroly nad tempem práce, konflikt práce a~péče, šikanu a~syndrom vyhoření. Národní akční plán pro duševní zdraví 2020--2030 řadí pracovní prostředí mezi oblasti prevence duševních obtíží. \ac{OECD} popisuje kvalitu práce jako vícerozměrnou kategorii, nikoli pouze jako vyjádření mezd a~pracovní doby. Pakliže má sociální dialog plnit adaptační funkci v~době digitalizace, musí se do~\aclgen{KV} dostat také právo na~odpojení, udržitelné výkonové normy a~ochrana před zhoršováním pracovních podmínek.~\cite{OECD_GoodJobsForAll_2018}\par

% \paragraph{Zdroje praktické části} propojují přehled poznání s~navazující empirickou analýzou. Kvantitativní část se opírá o~harmonizované databáze: evropské statistiky Eurostatu pro mzdy, příjmovou nerovnost a~daňový klín, srovnání pokrytí \acs{KS} a~odborové hustoty z~databází \acs{OECD} a~\acs{ETUI} a~českou mzdovou distribuci z~\acs{ISPV}, kterou lze členit na~pracoviště s~kolektivní smlouvou a~bez ní. Tyto zdroje umožňují zasadit \acacc{ČR} do~srovnání zemí \aclgen{EU} a~odlišit deklarovaný právní rámec od~jeho reálného dosahu.~\cite{OECD_UnionEmployerCoverage_2025}~\cite{etui_cba}\par

% Kvalitativní část navazuje studiem čtyř konkrétních smluv, které pokrývají různé konfigurace sociálního dialogu. \Ac{PKS} obchodního řetězce Penny Market je příkladem formalizovaného podnikové vyjednávání. Kolektivní smlouva podniku \ac{ČP} ukazuje silnou odborovou organizaci u~velkého zaměstnavatele s~problémy v~hospodaření. \ac{KSVS} v~zemědělství dokládá funkčnost rozšíření závaznosti podle \mbox{§~7}; a~odvětvová smlouva v~bankovnictví a~pojišťovnictví ilustruje sektorové vyjednávání s~nízkým pokrytím. Tyto smlouvy jsou ilustracemi, na~nichž lze ověřit, zda formální existence \acgen{KS} odpovídá reálné síle sociálního dialogu.~\cite{PennyMarket_KS2026}~\cite{KS_CeskaPosta_2025}~\cite{KSVS_OSPZV_ASO_2026}~\cite{KSVS_OSPPP_2026}\par
